% ==================================================================
%  chapters/preface.tex —— 原书前言 (p. I ~ II)
%  严格遵循《中文数学书籍 LaTeX 重排规范》：英文标点、忠实原文
% ==================================================================

\chapter*{前言}
\addcontentsline{toc}{chapter}{前言}
\markboth{前言}{前言}

数学分析是数学系最重要的基础课. 它对后继课程 (实变函数, 泛函分析, 拓扑, 微分几何) 与近代数学的学习与研究具有非常深远的影响和至关重要的作用. 一本优秀的数学分析教材必须包含传统微积分内容的精髓和分析能力与方法的传授, 也必须包含近代的内容, 其检验标准是若干年后能否涌现出一批高水准的应用数学人才和数学研究人才, 特别是一些数学顶尖人物. 作者从事数学分析教学几十年, 继承导师, 著名数学家吴文俊教授的一整套教学 (特别是教授数学分析的) 方法 (科大称之为``吴龙''), 并将其发扬光大, 因材施教, 在中国科技大学培养了一批国内外有名的数学家与数学工作者. 目前, 作者徐森林被特聘到华中师范大学数学与统计学学院, 并在数学试点班用此教材讲授数学分析, 效果显著.

本书的主要特色可归纳为以下几点.

\section*{1. 传统精髓内容的完善化}

书中包含了实数的各种引入, 七个实数连续性等价命题的论述; 给出了单变量与多变量的 Riemann 可积的各等价命题的证明; 讨论了微分中值定理, Taylor 公式余项的各种表达; 介绍了积分第一, 第二中值定理的描述, 隐函数存在性定理与反函数定理的两种不同的证法等内容.

\section*{2. 与后继课程的紧密结合, 使内容近代化}

本书在介绍经典微积分理论的同时, 将近代数学中许多重要概念, 理论恰到好处地引进分析教材中. 例如, 在积分理论中, 给出了 Lebesgue 定理: 函数 $f$ Riemann 可积的充要条件是 $f$ 几乎处处连续且有界; 详细讨论了 $\mathbb{R}^n$ 中的拓扑及相应的开集, 闭集, 聚点等概念, 描述了 $\mathbb{R}^n$ 中集合的紧致性, 连通性, 可数性, Hausdorff 性等拓扑不变性, 使读者站到拓扑的高度来理解零值定理, 介值定理, 最值定理与一致连续性定理. 引进外微分形式及外微分运算, 将经典 Newton-Leibniz 公式, 平面 Green 公式, 空间 Stokes 公式与 Gauss 公式统一为 Stokes 公式, 并对闭形式, 恰当形式与场论的对偶关系给出了全新的表述. 这不仅使教材内容本身近代化, 而且为学生在高年级学习拓扑, 实变函数, 泛函分析, 微分几何等课程提供了一个实际模型并打下良好的基础, 为经典数学与近代数学架设了一座桥梁.

\section*{3. 因材施教, 着重培养学生的研究与创新能力}

同一定理 (如零值定理, 一致连续性定理, Lagrange 中值定理, Cauchy 中值定理, 隐函数存在性定理与反函数定理等) 经常采用多种证法; 同一例题应用不同定理或不同方法解答, 这是本书又一特色. 它使学生广开思路, 积极锻炼思维能力, 使思维越来越敏捷与成熟. 书中举出大量例题是为了让读者得到一定的基本训练, 同时从定理的证明和典型实例的分析中掌握数学分析的技巧与方法. 习题共分三个层次: 练习题, 思考题与复习题. 练习题是基本题, 是为读者熟练掌握内容与方法设置的. 为提高学生对数学的浓厚兴趣及解题的能力, 设置了思考题. 为了让读者减少做题的障碍, 增强对数学的自信心, 其中有些题给出了提示. 实际上, 该节的标题就是最好的提示. 在每一章设置了大量复习题, 这些题不给提示, 因此大部分学生对它们会感到无从下手, 这些题是为少数想当数学家的学生特别设置的, 希望他们能深入思考, 自由发挥, 将它们一个一个地解答出来, 为将来的研究培养自己的创新能力. 如有困难, 我们还可撰写一本精练的学习指导书.

本书共分三册. 第一册内容包括数列极限, 函数极限与连续, 一元函数的导数与微分中值定理, Taylor 公式, 不定积分以及 Riemann 积分; 第二册内容包括 $\mathbb{R}^n$ 中的拓扑, $n$ 元函数的极限与连续, $n$ 元函数的微分学, 隐函数定理与反函数定理, $n$ 重积分, 第一型曲线, 曲面积分, 第二型曲线, 曲面积分, Stokes 定理, 外微分形式与场论; 第三册内容包括数项级数和各种收敛判别法, 函数项级数的一致收敛性及其性质, 含参变量反常积分的一致收敛性及其性质, Euler 积分 ($\Gamma$ 函数与 $B$ 函数), 幂级数与 Taylor 级数, Fourier 分析.

在写作本书的时候, 得到了华中师范大学数学与统计学学院领导和教师们的热情鼓励与大力支持, 作者们谨在此对他们表示诚挚的感谢. 博士生邓勤涛, 胡自胜, 薛琼, 硕士生金亚东, 鲍焱红等对本书的写作提出了许多宝贵意见, 使本书增色不少.

特别还要感谢的是清华大学出版社的曾刚, 刘颖, 王海燕, 他们为我们提供了本书出版的机会, 了却了我多年的心愿.

\vspace{1.5\baselineskip}
\begin{flushright}
  徐森林\\
  2005\,年\,6\,月于武汉
\end{flushright}
